% Pista ana-23j-q1 · Anàlisi
% Procedència: PAU Catalunya 2023 · Convocatòria de juny · Sèrie 1 · Qüestió 1
\documentclass[11pt,a4paper]{article}
\usepackage{pista}
\pistainfo{Catalunya 2023}{Convocatòria de juny}{Sèrie 1}

\begin{document}

\section*{\textcolor{blauPista}{Qüestió 1}}

\noindent Calculeu els coeficients $a$, $b$, $c$ i $d$ de la funció $f(x) = ax^{3} + bx^{2} + cx + d$ si sabem que l'equació de la recta tangent a la gràfica de la funció $f$ en el punt d'inflexió $(1, 0)$ és $y = -3x + 3$ i que la funció té un extrem relatiu en el punt de la gràfica d'abscissa $x = 0$. \hfill\textnormal{[2,5~punts]}

\pista{Cadascuna de les dades de l'enunciat es tradueix en una equació diferent. Per tant, al final tindràs un sistema d'equacions plantejat, que hauràs de resoldre. Para atenció amb la primera dada, perquè ``el punt $(1, 0)$ és d'inflexió'' amaga \emph{dues} condicions: que $(1, 0)$ pertany a la gràfica i que per $x=1$ s'anul·la $f''$. D'altra banda, la recta tangent a la gràfica en $(1, 0)$ té pendent $-3$, i ``extrem relatiu en $x = 0$'' també està donant una equació on intervé $f'(x)$. Si les comptes bé, obtindràs un sistema de quatre equacions amb les incògnites $a$, $b$, $c$ i $d$.}

\end{document}
